\documentclass[dvipdfmx,11pt]{beamer}

\usetheme{Madrid}
\usecolortheme{default}

\usepackage{amsmath,amssymb,bm}
\usepackage{listings}

\title{解釈可能な深層VARモデルの提案}
\subtitle{局所解釈性と効果的なスパース制約}
\author{栗田 煌矢}
\date{\today}

\begin{document}

\begin{frame}
  \titlepage
\end{frame}

\begin{frame}{問題意識}
  多変量時系列データから知見(依存関係や因果関係)を得る方法は様々である。
  
  おおざっぱに、既存モデルの特徴を挙げる。
  
  \begin{itemize}
    \item 線形モデル：複雑な非線形関係は捉えられない。検定や解釈が可能(例：VAR, VAR-LiNGAM)
    \item 非線形モデル：複雑な非線形関係を捉えうる。検定や解釈(例えば影響の正負)が不可能(例：DBN,Transfer Entropy)
    \item 両者共通：分析に使うラグを決める必要がある。
  \end{itemize}
  \vspace{0.5em}
  これらの短所をうまい具合に解決する、複雑でないモデルを作りたい。
\end{frame}

\begin{frame}{提案手法について}
  本手法は、多変量時系列データに対して
  解釈性を保ちながら深層学習手法を適用し、Granger因果構造(変数間の影響構造)を推定する。

  \vspace{0.5em}
  主な特徴は次の3点である。

  \begin{itemize}
    \item GVAR(Generalized Vector Autoregression)により、状態依存的かつ符号解釈が可能な自己回帰係数を学習する。
    \item Neural Granger Causality型の正則化により、回帰係数を適切にスパース化する。
    \item 学習可能な\texttt{causal\_gate}に近接勾配法を適用し、厳密なゼロを生成する。
  \end{itemize}
\end{frame}

\begin{frame}{背景：GVARとNeural Granger Causality}
  提案手法は、次の2つの既存手法を参考にしている。

  \begin{itemize}
    \item GVAR：状態依存的な係数行列をニューラルネットワークで生成する(Marcinkevičs, R., \& Vogt, J. E.,2021)。
    \item Neural-GC：ニューラルネットワークに構造的スパース性を導入し、Granger因果性を推定する(Tank et al.,2021)。
  \end{itemize}

  \vspace{0.5em}
  中心的な拡張は、
  \[
    \texttt{causal\_gate.shape} = [\text{lag}, \text{target}, \text{source}]
  \]
  という学習可能なゲートを導入した点である。
\end{frame}

\begin{frame}{GVAR(Marcinkevičs, R., \& Vogt, J. E.,2021)}
  $p$変量時系列を
  \[
    x_t = (x_{t,1}, \ldots, x_{t,p})^\top \in \mathbb{R}^{p}
  \]
  とする。自己回帰次数を$K$とすると、一般化VAR(GVAR)は次のように表される。

  \[
    x_t
    =
    \sum_{k=1}^{K}
    \Phi_k(x_{t-k}) x_{t-k}
    +
    \varepsilon_t
  \]

  ここで、$\varepsilon_t \in \mathbb{R}^{p}$は誤差項であり、
  $\Phi_k(x_{t-k}) \in \mathbb{R}^{p \times p}$はラグ$k$に対応する
  状態依存的な係数行列である。各成分で書くと、

  \[
    x_{t,i}
    =
    \sum_{k=1}^{K}
    \sum_{j=1}^{p}
    \Phi_{k,i,j}(x_{t-k}) x_{t-k,j}
    +
    \varepsilon_{t,i}
  \]

  $\Phi_{k,i,j}(x_{t-k})$はラグ$k$における変数$j$から
  変数$i$への局所的な影響を表す。
\end{frame}

\begin{frame}{Causal Gateによる構造制御(提案1)}
  本手法では、各ラグの係数行列 $\Phi_k(x_{t-k})$に対して
  静的なcausal gateを掛ける。

  \[
    G_k \in \mathbb{R}^{p \times p}
  \]

  有効係数行列は次で定義される。

  \[
    \tilde{\Phi}_k(x_{t-k})
    =
    G_k \odot \Phi_k(x_{t-k})
  \]

  \vspace{0.5em}
  ここで$\odot$はHadamard積である。
  動的な係数変動は$\Phi_k$が表現し、
  構造的な因果の有無を$G_k$が制御する。

  \vspace{0.5em}
  影響の「構造」は一定、影響の「度合い」は入力値依存という仮定を置いている。
\end{frame}

\begin{frame}{Granger非因果性の定義(Tank et al.,2021)}
  変数$x_j$が変数$x_i$をGranger causeしないための十分条件は、
  すべてのラグにおいて$x_j$から$x_i$への効果がゼロであることである。
  つまり、
  \[
    G_{1,i,j} = G_{2,i,j} = \cdots = G_{K,i,j} = 0
  \]

  causal gateの値が全てのラグにおいて0であれば、
  最終的な係数行列$\tilde{\Phi}_k(x_{t-k})$は、どんな入力値であろうと0になる。

  $\Phi_k(x_{t-k})$は入力値依存なので、同じラグでも0になる場合と
  ならない場合が考えられる。そのため、非因果性はcausal gateで定義した。
\end{frame}

\begin{frame}{予測}
  予測は、ラグごとの有効係数行列を用いて定義される。

  \[
    \hat{x}_t =
    \sum_{k=1}^{K}
    \tilde{\Phi}_k(x_{t-k}) x_{t-k}
  \]

  \vspace{0.5em}
  係数は入力依存のため、その形状は

  \[
    [\text{batch}, \text{lag}, \text{target}, \text{source}]
  \]

  であり、causal gateは入力依存でないため、

  \[
    [\text{lag}, \text{target}, \text{source}]
  \]

  の形を持つ。
\end{frame}

\begin{frame}{Temporal Smoothness(≒局所線形性)}
  GVARは時点ごとに状態依存係数を生成するため、
  係数が過度に振動しないように平滑化ペナルティを導入する。

  absolute modeでは次のペナルティを用いる。

  \[
    \mathcal{R}_{\mathrm{smooth}}
    =
    \frac{1}{|\mathcal{T}|}
    \sum_{t \in \mathcal{T}}
    \left\|
      \tilde{\Phi}_{t+1}
      -
      \tilde{\Phi}_t
    \right\|_F^2
  \]

 これにより、同じ変数でもラグをまたいで大きく係数の値(符号)が変わるという
 不自然な動きを抑制し、過学習防止や解釈の妥当性を実現する
\end{frame}

\begin{frame}{目的関数}
  学習目的関数は、予測損失、時間方向の平滑化、NGC正則化からなる。

  \[
    \mathcal{L}
    =
    \mathcal{L}_{\mathrm{pred}}
    +
    \lambda_{\mathrm{smooth}}
    \mathcal{R}_{\mathrm{smooth}}
    +
    \lambda_{\mathrm{ngc}}
    \mathcal{R}_{\mathrm{ngc}}
  \]

  予測損失は平均二乗誤差である。

  \[
    \mathcal{L}_{\mathrm{pred}}
    =
    \frac{1}{N}
    \sum_t
    \left\|
      x_t - \hat{x}_t
    \right\|_2^2
  \]

\end{frame}

\begin{frame}{Sparse Group Lasso(Tank et al.,2021)}
  本手法ではNGCのsparse group lassoをcausal gateに適用する。

  各target-source pair $(i,j)$について、ラグ方向のgate vectorを

  \[
    g_{i,j}
    =
    (G_{1,i,j}, \dots, G_{K,i,j})
  \]

  と定義する。

  正則化項は次である。

  \[
    \mathcal{R}_{\mathrm{ngc}}
    =
    w_{\mathrm{group}}
    \sum_{i,j}
    \|g_{i,j}\|_2
    +
    w_{\mathrm{l1}}
    \sum_{k,i,j}
    |G_{k,i,j}|
  \]

  group項はedge全体の削除、L1項はlag単位の削除に対応する。
\end{frame}

\begin{frame}{Iterative Shrinkage-Thresholding Algorithm(ISTA)}
  提案手法では、近接勾配法(ISTA)と呼ばれる閾値処理アルゴリズムにより、正則化
  で限りなく0に近づいた要素を厳密な0に更新する。

  まずsmooth partのみに対して勾配ステップを行う。

  \[
    \mathcal{L}_{\mathrm{smooth\text{-}part}}
    =
    \mathcal{L}_{\mathrm{pred}}
    +
    \lambda_{\mathrm{smooth}}
    \mathcal{R}_{\mathrm{smooth}}
  \]

  その後、causal gateに対して近接作用素を適用する。

  \[
    G^{(m+1)}
    =
    \mathrm{prox}_{
      \eta \lambda_{\mathrm{ngc}}
      \mathcal{R}_{\mathrm{ngc}}
    }
    \left(
      G^{(m+1/2)}
    \right)
  \]

  これにより、causal gateに厳密なゼロが生じる。
\end{frame}

\begin{frame}{近接更新の2段階}
  sparse group lassoの近接更新は2段階で行われる。

  \vspace{0.5em}
  まず、elementwise soft-thresholdingによりlag単位の厳密なスパース性を作る。

  \[
    G_{k,i,j}
    \leftarrow
    \mathrm{sign}(G_{k,i,j})
    \left(
      |G_{k,i,j}|
      -
      \eta \lambda_{\mathrm{ngc}} w_{\mathrm{l1}}
    \right)_+
  \]

  次に、各edge $(i,j)$のlag vectorに対してgroup soft-thresholdingを行う。

  \[
    g_{i,j}
    \leftarrow
    \left(
      1 -
      \frac{
        \eta \lambda_{\mathrm{ngc}} w_{\mathrm{group}}
      }{
        \|g_{i,j}\|_2
      }
    \right)_+
    g_{i,j}
  \]
\end{frame}

\begin{frame}{SGD・Adamと近接勾配法(ISTA)の違い}
  XNeural VARでは、\texttt{adam}と\texttt{ista}の2種類のoptimizerを選択できる。

  \begin{itemize}
    \item \texttt{SGD・Adam}
      \begin{itemize}
        \item 目的関数全体を通常の勾配降下法ベースで最適化する。
        \item 実装は簡単だが、目的の要素(ここではcausal gate)が厳密なゼ0になりにくい。
      \end{itemize}

    \item \texttt{ista}
      \begin{itemize}
        \item 通常の勾配法を行った後、gateに近接作用素を適用する。
        \item 厳密な0を作るため、Granger構造推定に適している。
      \end{itemize}
  \end{itemize}
\end{frame}

\begin{frame}{Graph Inference}
  学習後、モデルは全データに対して有効係数テンソルを計算する。

  係数ベースのcausal strengthは、係数の絶対値をbatch方向およびlag方向に集約して得る。

  \[
    S_{i,j}
    =
    \max_{t,k}
    \left|
      \tilde{\Phi}_{k,i,j}(x_{t-k})
    \right|
  \]

  ただし、最終的なcausal graphはcoefficient thresholdingではなく、
  gate normに基づいて推定する。

  \[
    \hat{A}_{i,j}
    =
    \mathbf{1}
    \left[
      \|g_{i,j}\|_2 > \tau
    \right]
  \]
\end{frame}

\begin{frame}{提案手法のパッケージ化(https://github.com/koyarctan/Xneural-VAR)}
  XNeural VARは、機能ごとにモジュール化されている。

  \begin{itemize}
    \item \texttt{xneural\_var.models}\\
    GVAR/SENNモデルとcausal gateを実装する。
    \item \texttt{xneural\_var.regularizers}\\
    sparse group lasso、hierarchical group lasso、近接作用素を実装する。
    \item \texttt{xneural\_var.training}\\
    学習ループ、ログ出力、結果オブジェクトを扱う。
  \end{itemize}
\end{frame}

\begin{frame}{まとめ}
  本手法の要点は次の通りである。

  \begin{itemize}
    \item GVARにより、非線形かつ状態依存的な自己回帰ダイナミクスを表現する。
    \item causal gateにより、動的係数と静的なGranger構造を分離する。
    \item sparse group lassoにより、edge全体とlag単位の両方をスパース化する。
    \item ISTAにより、causal gateに厳密なゼロを生成する。
    \item 平滑化制約によって自然な状態依存係数の推移を実現。解釈の妥当性を生む。
  \end{itemize}
\end{frame}

\end{document}